\chapter{Judges and Officials Rules}

\section{Basketball Director}

The Basketball Director is the head organizer and administrator of basketball competitions.
The Basketball Director is responsible for the logistics and equipment for the basketball
competitions. The Basketball Director is in charge of deciding and specifying the
tournament structure, and of keeping events running on schedule. The Basketball
Director is responsible for resolving the general protests not concerning specific games,
except for the choice of referees for a game.

The Basketball Director may be assisted by a Basketball Co-Director. They shall not be
players of the same basketball team in the competition. If a protest concerns the team of
one of the two Directors, the final decision on the protest is made by the other Director.
In cases with no conflicts of interest but in which they disagree, the final decision is
made by the (primary) Director.

\section{Game Officials}

As with Rule Eight of the FIBA Rulebook, game officials shall consist of two or three referees 
and at least two table officials (at least the scorer and a timer/assistant scorer). 
Referees shall have all of the responsibilities and powers described in Articles 45, 46, and 47 
except for articles covering the referees uniforms. The referees are responsible for resolving protests concerning the game.

Referees should be selected for the game based on the following criteria:

\begin{itemize}
  \item Referees should not be teammates or members of the clubs of either team (even members not playing in the tournament)
  \item Referees should not include only players/participants from one team within the same tournament
  \item When the tournament consists of teams from more than one region, referees should not include only players/participants from one region
  \item Referees should not include players in the same tournament
\end{itemize}

If the referee pool is insufficient the directors may relax these restrictions starting at the bottom 
of the list and working upward. The directors may choose to limit the pool based on technical referee skills.

At Unicon, for all playoff games in the A tournament, experienced referees who are not playing in the A tournament 
and not associated with any teams in the A tournament should be used. These would typically be certified on-foot basketball 
referees paid via the Unicon budget. One such referee per game should also be planned for other important games preceding 
the playoffs (this would depend on the tournament format), allowing them to get used to the unicycle basketball particularities before the playoffs.

\newpage
\section{Referees' Signals}
Referees should use the hand signals in Appendix A of the FIBA rules. In addition, the following hand signals may be used.

\begin{itemize}
  \item {\underline{Unmounted play:}} Straight arm, with the hand in front of the hips. The palm is facing the floor and simulates a sweeping (wrist move).
  \item {\underline{Rolling the ball on floor:}} One vertically rotating fist, below the waist.
  \item {\underline{Ball pick-up interference:}} Both arms parallel pointing down to the side, with hand palms facing each other.
\end{itemize}

Referees should try to always use the appropriate signals, especially in international competitions. 
In any case, referees should at least use the appropriate signals for interrupting the game 
(one arm raised, with open palm for a violation or with a clenched fist for a foul), and then for indicating the direction
of play (point in direction of play, arm parallel to sidelines).